\begin{explanationbox}
	\textbf{\textcolor{CExplanation}{一句话直觉}}

	长方体像一个盒子，外接球刚好套住它，球的直径正好是盒子的体对角线。

	\medskip
	\textbf{\textcolor{CExplanation}{核心拆解}}
	\begin{description}[style=nextline, leftmargin=2.8em, labelsep=0.5em, itemsep=0.45em]
		\item[\textbf{要点一（球心位置）：}] 长方体外接球的球心在长方体中心，也就是任意一条体对角线的中点。
		\item[\textbf{要点二（体对角线）：}] 体对角线长度为 $\sqrt{a^2+b^2+c^2}$，可由两次勾股定理或空间两点间距离公式得到。
		\item[\textbf{要点三（半径关系）：}] 体对角线的两个端点都在球面上，并且体对角线经过球心，所以体对角线是外接球直径，半径是它的一半。
	\end{description}

	\medskip
	\textbf{\textcolor{CExplanation}{几何本质（体对角线）}}

	长方体的任意一条体对角线都穿过长方体中心，且两端点都在外接球球面上。因此，体对角线就是外接球的一条直径。

	\medskip
	\textbf{\textcolor{CExplanation}{代数意义（平方和开方）}}

	公式 $R=\frac{1}{2}\sqrt{a^2+b^2+c^2}$ 体现了长方体三个互相垂直方向长度的平方和再开方，与空间距离公式一致。

	\medskip
	\textbf{\textcolor{CExplanation}{考点价值（高频公式）}}
	\begin{itemize}[leftmargin=*, itemsep=0.5em, topsep=0.4em]
		\item \textbf{考法一（直接套用）：} 已知长方体长、宽、高，直接代入公式求外接球半径。
		\item \textbf{考法二（反求边长）：} 已知外接球半径，结合长、宽、高之间的关系列方程。
		\item \textbf{考法三（补形转化）：} 若题目中的几何体可看作某个长方体的一部分顶点，可转化为长方体外接球问题。
	\end{itemize}

	\medskip
	\textbf{\textcolor{CExplanation}{顿悟点}}

	\textbf{体对角线就是直径}，所以只要能找到长方体的长、宽、高，就能得到外接球半径。

	\medskip
	\textbf{\textcolor{CExplanation}{使用场景}}
	\begin{itemize}[leftmargin=*, itemsep=0.5em, topsep=0.4em]
		\item \textbf{场景一（长方体直接计算）：} 在立体几何题中，遇到长方体外接球，直接套用公式。
		\item \textbf{场景二（球与多面体补形）：} 若原几何体的顶点可看作某个长方体的部分顶点，且这些顶点确定的外接球与该长方体外接球相同，则可应用此公式。
	\end{itemize}
\end{explanationbox}